% assets/a3-geometry.svg — 아지무스 추진기 둘의 기하와 배분행렬 B
%
%   bash _tools/tikz2svg.sh assets/src/a3-geometry.tex assets/a3-geometry.svg
%
% 이 그림이 말하는 것
%   추진기가 돌아가면 한 추진기가 내놓는 것이 **힘 하나가 아니라 성분 둘**이 된다.
%   6주차의 추진기 하나 = 미지수 하나였던 자리가, 여기서는 미지수 둘이 된다.
%       좌현 (F_x,L, F_y,L)   우현 (F_x,R, F_y,R)   -> 미지수 넷, 식 셋
%   배분행렬 B 의 네 열은 그 넷이 X, Y, N 에 각각 얼마나 들어가는지를 적은 것이다.
%
% 좌표 — 배를 위에서 본다. 선체 좌표계
%   x_b = 선수 방향(그림 위), y_b = 우현 방향(그림 오른쪽)
%   \BP{xb}{yb} 로 적는다. 1 m = 20 mm
%   화면 x = y_b, 화면 y = x_b 이므로 매크로 안에서 인자 순서가 뒤집힌다
%   \BQ 는 괄호 없는 짝 — ([shift={...}]\BQ{..}{..}) 로 쓴다
%
% 각도 부호
%   델타는 x_b 에서 y_b 쪽으로 재는 각 — 화면에서는 위에서 오른쪽으로(시계방향)
%   좌현 +40 도 (길이 2.2 m), 우현 -25 도 (길이 1.2 m) 로 그렸다.
%   둘이 서로 다른 값을 가질 수 있다는 것이 요점이므로 일부러 다르게 그린다
%
% 겹침을 피하려고 정한 것
%   두 추력 벡터의 끝을 세로로 12 mm 떼어 놓았다 (길이를 2.2 m 와 1.2 m 로 다르게)
%   F_y 두 점선이 같은 높이에 놓이면 라벨이 겹친다 — 그래서 길이를 다르게 잡았다
%   델타 라벨은 호의 이등분선 위 반지름 11 mm 자리에 둔다

\documentclass[border=6pt]{standalone}
% capstone-style.tex — 이 볼트의 TikZ 그림이 공유하는 스타일
%
% 쓰는 법:  % capstone-style.tex — 이 볼트의 TikZ 그림이 공유하는 스타일
%
% 쓰는 법:  % capstone-style.tex — 이 볼트의 TikZ 그림이 공유하는 스타일
%
% 쓰는 법:  \input{capstone-style}   (assets/src/ 안에서)
% 빌드:     bash _tools/tikz2svg.sh assets/src/w02-node-topic.tex assets/w02-node-topic.svg
%
% 왜 TikZ 인가
%   손으로 SVG 좌표를 쓰면 화살촉·선 굵기·글자 크기가 그림마다 미묘하게 달라진다.
%   그것이 "자료가 어설퍼 보인다"의 정체다. 스타일을 한 번만 정하고 상대좌표로
%   배치하면 좌표를 고쳐도 그림이 무너지지 않는다.
%   수식은 본문과 같은 TeX 글꼴로 나온다.
%
% 한글
%   ko.TeX(DVI 경로)를 쓴다. XeLaTeX 은 TikZ 그래픽을 PDF special 로 내보내는데
%   dvisvgm 이 그것을 받으려면 Ghostscript < 10.01 이 필요하다. 이 PC 에는
%   10.07 이 깔려 있어 그 경로가 막힌다. latex -> DVI -> dvisvgm 이 유일하게
%   동작하는 조합이다.

\usepackage[hangul]{kotex}
\usepackage{tikz}
\usepackage{amsmath}
\usepackage{xcolor}
\usetikzlibrary{arrows.meta, positioning, calc, fit, backgrounds, decorations.pathmorphing}

% ---- 색 --------------------------------------------------------------------
% 강조는 하나만. 나머지는 검정. 흑백 인쇄에서도 읽혀야 한다.
\definecolor{ink}{HTML}{1A1A1A}
\definecolor{soft}{HTML}{666666}
\definecolor{accent}{HTML}{7C3AED}
\definecolor{warn}{HTML}{B91C1C}

% 계층 색 — Simulink 모델의 블록 색과 같은 규칙을 쓴다
\definecolor{guid}{HTML}{1D4ED8}    % 유도
\definecolor{ctrl}{HTML}{EA580C}    % 제어
\definecolor{thr}{HTML}{D97706}     % 추진기
\definecolor{plant}{HTML}{16A34A}   % 운동모델
\definecolor{nav}{HTML}{7C3AED}     % 항법 · 신호처리
\definecolor{pathc}{HTML}{0D9488}   % 계획 경로 · 항적
\definecolor{errc}{HTML}{C2410C}    % 오차
\definecolor{flow}{HTML}{0369A1}    % 조류 · 바람

% ---- 선과 화살촉 -----------------------------------------------------------
\tikzset{
  sig/.style   = {ink, line width=0.5pt, -{Stealth[length=4pt, width=3pt]}},
  wire/.style  = {ink, line width=0.5pt},
  fb/.style    = {ink, line width=0.5pt, densely dashed,
                  -{Stealth[length=4pt, width=3pt]}},
  asig/.style  = {accent, line width=0.5pt, densely dashed,
                  -{Stealth[length=4pt, width=3pt]}},
  % 블록 — 흰 바탕, 직각, 검은 테두리. 색은 테두리에만 준다
  blk/.style   = {draw=ink, line width=0.55pt, fill=white,
                  minimum width=18mm, minimum height=8.5mm, inner sep=2pt, font=\small},
  gblk/.style  = {blk, draw=guid},
  cblk/.style  = {blk, draw=ctrl},
  tblk/.style  = {blk, draw=thr},
  pblk/.style  = {blk, draw=plant},
  nblk/.style  = {blk, draw=nav},
  % 합산점 · 분기점
  sum/.style   = {draw=ink, line width=0.55pt, fill=white, circle,
                  inner sep=0pt, minimum size=4.6mm, font=\scriptsize},
  dot/.style   = {ink, fill=ink, circle, inner sep=0pt, minimum size=1.5mm},
  % 글자
  sname/.style = {font=\small, inner sep=1.5pt},
  note/.style  = {font=\scriptsize, soft, inner sep=1.5pt},
  wnote/.style = {font=\scriptsize, warn, inner sep=1.5pt},
  anote/.style = {font=\scriptsize, accent, inner sep=1.5pt},
  axis/.style  = {ink, line width=0.5pt},
  tick/.style  = {font=\scriptsize, soft},
  % 묶음 상자 — 서브시스템 표시
  grp/.style   = {draw=soft, line width=0.4pt, densely dotted, rounded corners=1mm,
                  inner sep=3mm},
}
   (assets/src/ 안에서)
% 빌드:     bash _tools/tikz2svg.sh assets/src/w02-node-topic.tex assets/w02-node-topic.svg
%
% 왜 TikZ 인가
%   손으로 SVG 좌표를 쓰면 화살촉·선 굵기·글자 크기가 그림마다 미묘하게 달라진다.
%   그것이 "자료가 어설퍼 보인다"의 정체다. 스타일을 한 번만 정하고 상대좌표로
%   배치하면 좌표를 고쳐도 그림이 무너지지 않는다.
%   수식은 본문과 같은 TeX 글꼴로 나온다.
%
% 한글
%   ko.TeX(DVI 경로)를 쓴다. XeLaTeX 은 TikZ 그래픽을 PDF special 로 내보내는데
%   dvisvgm 이 그것을 받으려면 Ghostscript < 10.01 이 필요하다. 이 PC 에는
%   10.07 이 깔려 있어 그 경로가 막힌다. latex -> DVI -> dvisvgm 이 유일하게
%   동작하는 조합이다.

\usepackage[hangul]{kotex}
\usepackage{tikz}
\usepackage{amsmath}
\usepackage{xcolor}
\usetikzlibrary{arrows.meta, positioning, calc, fit, backgrounds, decorations.pathmorphing}

% ---- 색 --------------------------------------------------------------------
% 강조는 하나만. 나머지는 검정. 흑백 인쇄에서도 읽혀야 한다.
\definecolor{ink}{HTML}{1A1A1A}
\definecolor{soft}{HTML}{666666}
\definecolor{accent}{HTML}{7C3AED}
\definecolor{warn}{HTML}{B91C1C}

% 계층 색 — Simulink 모델의 블록 색과 같은 규칙을 쓴다
\definecolor{guid}{HTML}{1D4ED8}    % 유도
\definecolor{ctrl}{HTML}{EA580C}    % 제어
\definecolor{thr}{HTML}{D97706}     % 추진기
\definecolor{plant}{HTML}{16A34A}   % 운동모델
\definecolor{nav}{HTML}{7C3AED}     % 항법 · 신호처리
\definecolor{pathc}{HTML}{0D9488}   % 계획 경로 · 항적
\definecolor{errc}{HTML}{C2410C}    % 오차
\definecolor{flow}{HTML}{0369A1}    % 조류 · 바람

% ---- 선과 화살촉 -----------------------------------------------------------
\tikzset{
  sig/.style   = {ink, line width=0.5pt, -{Stealth[length=4pt, width=3pt]}},
  wire/.style  = {ink, line width=0.5pt},
  fb/.style    = {ink, line width=0.5pt, densely dashed,
                  -{Stealth[length=4pt, width=3pt]}},
  asig/.style  = {accent, line width=0.5pt, densely dashed,
                  -{Stealth[length=4pt, width=3pt]}},
  % 블록 — 흰 바탕, 직각, 검은 테두리. 색은 테두리에만 준다
  blk/.style   = {draw=ink, line width=0.55pt, fill=white,
                  minimum width=18mm, minimum height=8.5mm, inner sep=2pt, font=\small},
  gblk/.style  = {blk, draw=guid},
  cblk/.style  = {blk, draw=ctrl},
  tblk/.style  = {blk, draw=thr},
  pblk/.style  = {blk, draw=plant},
  nblk/.style  = {blk, draw=nav},
  % 합산점 · 분기점
  sum/.style   = {draw=ink, line width=0.55pt, fill=white, circle,
                  inner sep=0pt, minimum size=4.6mm, font=\scriptsize},
  dot/.style   = {ink, fill=ink, circle, inner sep=0pt, minimum size=1.5mm},
  % 글자
  sname/.style = {font=\small, inner sep=1.5pt},
  note/.style  = {font=\scriptsize, soft, inner sep=1.5pt},
  wnote/.style = {font=\scriptsize, warn, inner sep=1.5pt},
  anote/.style = {font=\scriptsize, accent, inner sep=1.5pt},
  axis/.style  = {ink, line width=0.5pt},
  tick/.style  = {font=\scriptsize, soft},
  % 묶음 상자 — 서브시스템 표시
  grp/.style   = {draw=soft, line width=0.4pt, densely dotted, rounded corners=1mm,
                  inner sep=3mm},
}
   (assets/src/ 안에서)
% 빌드:     bash _tools/tikz2svg.sh assets/src/w02-node-topic.tex assets/w02-node-topic.svg
%
% 왜 TikZ 인가
%   손으로 SVG 좌표를 쓰면 화살촉·선 굵기·글자 크기가 그림마다 미묘하게 달라진다.
%   그것이 "자료가 어설퍼 보인다"의 정체다. 스타일을 한 번만 정하고 상대좌표로
%   배치하면 좌표를 고쳐도 그림이 무너지지 않는다.
%   수식은 본문과 같은 TeX 글꼴로 나온다.
%
% 한글
%   ko.TeX(DVI 경로)를 쓴다. XeLaTeX 은 TikZ 그래픽을 PDF special 로 내보내는데
%   dvisvgm 이 그것을 받으려면 Ghostscript < 10.01 이 필요하다. 이 PC 에는
%   10.07 이 깔려 있어 그 경로가 막힌다. latex -> DVI -> dvisvgm 이 유일하게
%   동작하는 조합이다.

\usepackage[hangul]{kotex}
\usepackage{tikz}
\usepackage{amsmath}
\usepackage{xcolor}
\usetikzlibrary{arrows.meta, positioning, calc, fit, backgrounds, decorations.pathmorphing}

% ---- 색 --------------------------------------------------------------------
% 강조는 하나만. 나머지는 검정. 흑백 인쇄에서도 읽혀야 한다.
\definecolor{ink}{HTML}{1A1A1A}
\definecolor{soft}{HTML}{666666}
\definecolor{accent}{HTML}{7C3AED}
\definecolor{warn}{HTML}{B91C1C}

% 계층 색 — Simulink 모델의 블록 색과 같은 규칙을 쓴다
\definecolor{guid}{HTML}{1D4ED8}    % 유도
\definecolor{ctrl}{HTML}{EA580C}    % 제어
\definecolor{thr}{HTML}{D97706}     % 추진기
\definecolor{plant}{HTML}{16A34A}   % 운동모델
\definecolor{nav}{HTML}{7C3AED}     % 항법 · 신호처리
\definecolor{pathc}{HTML}{0D9488}   % 계획 경로 · 항적
\definecolor{errc}{HTML}{C2410C}    % 오차
\definecolor{flow}{HTML}{0369A1}    % 조류 · 바람

% ---- 선과 화살촉 -----------------------------------------------------------
\tikzset{
  sig/.style   = {ink, line width=0.5pt, -{Stealth[length=4pt, width=3pt]}},
  wire/.style  = {ink, line width=0.5pt},
  fb/.style    = {ink, line width=0.5pt, densely dashed,
                  -{Stealth[length=4pt, width=3pt]}},
  asig/.style  = {accent, line width=0.5pt, densely dashed,
                  -{Stealth[length=4pt, width=3pt]}},
  % 블록 — 흰 바탕, 직각, 검은 테두리. 색은 테두리에만 준다
  blk/.style   = {draw=ink, line width=0.55pt, fill=white,
                  minimum width=18mm, minimum height=8.5mm, inner sep=2pt, font=\small},
  gblk/.style  = {blk, draw=guid},
  cblk/.style  = {blk, draw=ctrl},
  tblk/.style  = {blk, draw=thr},
  pblk/.style  = {blk, draw=plant},
  nblk/.style  = {blk, draw=nav},
  % 합산점 · 분기점
  sum/.style   = {draw=ink, line width=0.55pt, fill=white, circle,
                  inner sep=0pt, minimum size=4.6mm, font=\scriptsize},
  dot/.style   = {ink, fill=ink, circle, inner sep=0pt, minimum size=1.5mm},
  % 글자
  sname/.style = {font=\small, inner sep=1.5pt},
  note/.style  = {font=\scriptsize, soft, inner sep=1.5pt},
  wnote/.style = {font=\scriptsize, warn, inner sep=1.5pt},
  anote/.style = {font=\scriptsize, accent, inner sep=1.5pt},
  axis/.style  = {ink, line width=0.5pt},
  tick/.style  = {font=\scriptsize, soft},
  % 묶음 상자 — 서브시스템 표시
  grp/.style   = {draw=soft, line width=0.4pt, densely dotted, rounded corners=1mm,
                  inner sep=3mm},
}


\begin{document}
\begin{tikzpicture}[x=1mm, y=1mm]

\pgfmathsetmacro{\BS}{20}                        % 1 m = 20 mm
\newcommand{\BQ}[2]{{(#2)*\BS},{(#1)*\BS}}       % 괄호 없는 짝
\newcommand{\BP}[2]{(\BQ{#1}{#2})}               % #1 = x_b (선수), #2 = y_b (우현)

\pgfmathsetmacro{\bb}{1.027135}                  % half_beam [m]
\pgfmathsetmacro{\LL}{2.373776}                  % l_thr     [m]

% ===========================================================================
% (가) 추진기가 앉은 자리와 돌아간 각도
% ===========================================================================
\begin{scope}[shift={(0,0)}]

  \node[font=\small\bfseries, anchor=north west] at (-34,62)
       {(가) 추진기 둘이 돌아가면 미지수가 넷이 된다};

  % --- 선체 외곽선 -------------------------------------------------------
  \draw[ink, line width=0.6pt, fill=white]
       \BP{-2.75}{-\bb} -- \BP{1.3}{-\bb} -- \BP{2.4}{0} -- \BP{1.3}{\bb}
       -- \BP{-2.75}{\bb} -- cycle;
  \draw[soft, line width=0.3pt, densely dashed] \BP{-2.85}{0} -- \BP{2.5}{0};

  % --- 선체 좌표축 -------------------------------------------------------
  \draw[nav, line width=0.6pt, -{Stealth[length=3.4pt,width=2.4pt]}]
       \BP{0}{0} -- \BP{2.1}{0};
  \node[font=\scriptsize, nav, anchor=south east] at \BP{2.12}{-0.08} {$x_b$ (선수)};
  \draw[nav, line width=0.6pt, -{Stealth[length=3.4pt,width=2.4pt]}]
       \BP{0}{0} -- \BP{0}{0.78};
  \node[font=\scriptsize, nav, anchor=north east] at \BP{-0.08}{0.82} {$y_b$ (우현)};
  \node[dot] at \BP{0}{0} {};
  \node[font=\scriptsize, soft, anchor=north east] at \BP{-0.12}{-0.08} {무게중심};

  % --- 좌현 추진기 — 델타 = +40 도, 길이 2.2 m --------------------------
  %     끝점 x_b = -0.6885,  y_b = +0.3871
  \draw[guid, line width=0.4pt, densely dashed]
       \BP{-\LL}{-\bb} -- \BP{-0.6885}{-\bb};
  \draw[guid, line width=0.4pt, densely dashed]
       \BP{-0.6885}{-\bb} -- \BP{-0.6885}{0.3871};
  \node[font=\scriptsize, guid, anchor=east]  at \BP{-1.50}{-1.12} {$F_{x,L}$};
  \node[font=\scriptsize, guid, anchor=south] at \BP{-0.6885}{-0.32} {$F_{y,L}$};
  \draw[thr, line width=1.3pt, -{Stealth[length=4.4pt,width=3.2pt]}]
       \BP{-\LL}{-\bb} -- \BP{-0.6885}{0.3871};
  \node[font=\scriptsize, thr, anchor=west] at \BP{-0.62}{0.46} {$T_L$};
  \draw[ctrl, line width=0.5pt, -{Stealth[length=2.6pt,width=2pt]}]
       ([shift={(0,15)}]\BQ{-\LL}{-\bb}) arc (90:50:15mm);
  \node[font=\scriptsize, ctrl, anchor=west]
       at ([shift={(3.76,10.34)}]\BQ{-\LL}{-\bb}) {$\delta_L$};
  \node[circle, draw=thr, fill=white, line width=0.7pt, minimum size=3.4mm,
        inner sep=0pt] at \BP{-\LL}{-\bb} {};

  % --- 우현 추진기 — 델타 = -25 도, 길이 1.2 m --------------------------
  %     끝점 x_b = -1.2862,  y_b = +0.5200
  \draw[ctrl, line width=0.4pt, densely dashed]
       \BP{-\LL}{\bb} -- \BP{-1.2862}{\bb};
  \draw[ctrl, line width=0.4pt, densely dashed]
       \BP{-1.2862}{\bb} -- \BP{-1.2862}{0.5200};
  \node[font=\scriptsize, ctrl, anchor=west]  at \BP{-1.80}{1.12} {$F_{x,R}$};
  \node[font=\scriptsize, ctrl, anchor=north] at \BP{-1.33}{0.78} {$F_{y,R}$};
  \draw[thr, line width=1.3pt, -{Stealth[length=4.4pt,width=3.2pt]}]
       \BP{-\LL}{\bb} -- \BP{-1.2862}{0.5200};
  \node[font=\scriptsize, thr, anchor=south east] at \BP{-1.24}{0.56} {$T_R$};
  \draw[ctrl, line width=0.5pt, -{Stealth[length=2.6pt,width=2pt]}]
       ([shift={(0,15)}]\BQ{-\LL}{\bb}) arc (90:115:15mm);
  \node[font=\scriptsize, ctrl, anchor=east]
       at ([shift={(-2.39,10.74)}]\BQ{-\LL}{\bb}) {$\delta_R$};
  \node[circle, draw=thr, fill=white, line width=0.7pt, minimum size=3.4mm,
        inner sep=0pt] at \BP{-\LL}{\bb} {};

  % --- 모멘트 팔 b — 선체 아래 바깥 --------------------------------------
  \draw[errc, line width=0.45pt, {Stealth[length=2.6pt]}-{Stealth[length=2.6pt]}]
       \BP{-3.10}{-\bb} -- \BP{-3.10}{0};
  \node[font=\scriptsize, errc, anchor=north] at \BP{-3.20}{-0.52} {$b = 1.027$ m};
  \draw[errc, line width=0.25pt, densely dotted] \BP{-3.20}{-\bb} -- \BP{-2.70}{-\bb};
  \draw[errc, line width=0.25pt, densely dotted] \BP{-3.20}{0}    -- \BP{-2.80}{0};

  % --- 모멘트 팔 L — 선체 오른쪽 바깥 ------------------------------------
  \draw[errc, line width=0.45pt, {Stealth[length=2.6pt]}-{Stealth[length=2.6pt]}]
       \BP{0}{1.70} -- \BP{-\LL}{1.70};
  \node[font=\scriptsize, errc, anchor=west] at \BP{-1.19}{1.78} {$L = 2.374$ m};
  \draw[errc, line width=0.25pt, densely dotted] \BP{0}{1.78}    -- \BP{0}{1.10};
  \draw[errc, line width=0.25pt, densely dotted] \BP{-\LL}{1.78} -- \BP{-\LL}{1.10};

  % --- 아래 설명 ---------------------------------------------------------
  \node[anchor=north west, font=\scriptsize, text width=90mm] at (-34,-76)
       {\textbf{돌아가는 추진기 하나는 미지수 둘이다} \\[4pt]
        추진기가 $\delta$ 만큼 돌아가 추력 $T$ 를 내면, 그 힘은 선체 앞뒤로 미는 몫과
        옆으로 미는 몫으로 갈린다 \\[4pt]
        \hspace*{6mm}$F_x = T\cos\delta$, \qquad $F_y = T\sin\delta$ \\[4pt]
        5 $\cdot$ 6주차의 고정 추진기는 $\delta \equiv 0$ 이었다 — 그래서 $F_y$ 가
        아예 없었고, 추진기 하나가 미지수 하나였다 \\[4pt]
        $\;\to\;$ 각도가 생기면서 미지수가 \textbf{둘에서 넷}으로 늘었다.
        식은 그대로 $X, Y, N$ 셋이다};
\end{scope}

% ===========================================================================
% (나) 배분행렬 B — 위 그림을 그대로 옮긴 것
% ===========================================================================
\begin{scope}[shift={(78,0)}]

  \node[font=\small\bfseries, anchor=north west] at (0,62)
       {(나) 배분행렬 $\mathbf{B}$ 의 네 열은 (가)를 옮겨 적은 것이다};

  \node[anchor=north, font=\small] at (46,50)
       {$\begin{bmatrix} X \\ Y \\ N \end{bmatrix}
         = \underbrace{\begin{bmatrix}
             \textcolor{guid}{1} & \textcolor{guid}{0}  & \textcolor{ctrl}{1}  & \textcolor{ctrl}{0} \\
             \textcolor{guid}{0} & \textcolor{guid}{1}  & \textcolor{ctrl}{0}  & \textcolor{ctrl}{1} \\
             \textcolor{guid}{b} & \textcolor{guid}{-L} & \textcolor{ctrl}{-b} & \textcolor{ctrl}{-L}
           \end{bmatrix}}_{\textstyle \mathbf{B}\ (3\times 4)}
           \underbrace{\begin{bmatrix}
             \textcolor{guid}{F_{x,L}} \\ \textcolor{guid}{F_{y,L}} \\
             \textcolor{ctrl}{F_{x,R}} \\ \textcolor{ctrl}{F_{y,R}}
           \end{bmatrix}}_{\textstyle \mathbf{u}_e}$};

  \node[anchor=north, font=\scriptsize] at (46,20)
       {\textcolor{guid}{파란 두 열 = 좌현} \qquad \textcolor{ctrl}{주황 두 열 = 우현}};

  \draw[soft, line width=0.3pt] (0,14) -- (92,14);

  \node[anchor=north west, font=\scriptsize, text width=92mm] at (0,9)
       {\textbf{세 줄을 읽는 법} \\[4pt]
        첫째 줄 — 앞뒤로 미는 몫. 네 성분 가운데 $x$ 성분 둘이 그대로 더해진다 \\[2pt]
        \hspace*{6mm}$X = F_{x,L} + F_{x,R}$ \\[4pt]
        둘째 줄 — 옆으로 미는 몫. $y$ 성분 둘이 그대로 더해진다 \\[2pt]
        \hspace*{6mm}$Y = F_{y,L} + F_{y,R}$ \\[4pt]
        셋째 줄 — 돌리는 몫. 추진기가 $(x_r, y_r)$ 에 앉아 있으면 모멘트는
        $N = x_r F_y - y_r F_x$ 한 줄로 나온다 \\[3pt]
        \hspace*{6mm}좌현 $(-L,\, -b)$ \; $\to$ \;
        $\textcolor{guid}{+\,b\,F_{x,L} - L\,F_{y,L}}$ \\[2pt]
        \hspace*{6mm}우현 $(-L,\, +b)$ \; $\to$ \;
        $\textcolor{ctrl}{-\,b\,F_{x,R} - L\,F_{y,R}}$ \\[4pt]
        $\;\to\;$ $\mathbf{B}$ 에는 \textbf{기하만} 들어 있다. 각도는 한 칸도 없다};

  \draw[soft, line width=0.3pt] (0,-50) -- (92,-50);

  \node[anchor=north west, font=\scriptsize, text width=92mm] at (0,-55)
       {\textbf{성분 넷을 추진기 지령 둘로 되돌린다} — 극좌표 변환 \\[4pt]
        \hspace*{6mm}$T = \sqrt{F_x^2 + F_y^2}$, \qquad
        $\delta = \operatorname{atan2}(F_y,\, F_x)$ \\[4pt]
        배분이 내놓는 것은 $\mathbf{u}_e$ 네 수이고, 서보와 추진기가 받는 것은
        각도 둘 $\cdot$ 추력 둘이다. 그 사이를 잇는 것이 이 두 줄뿐이다};

  \node[anchor=north west, font=\scriptsize, warn, text width=92mm] at (0,-81)
       {미지수 넷 $\cdot$ 식 셋 — \textbf{해가 하나로 정해지지 않는다}.
        같은 $X, Y, N$ 을 내는 $\mathbf{u}_e$ 가 직선 하나를 이룬다.
        그 직선에서 어느 점을 고를지가 배분 문제다 (1-4 $\cdot$ 1-7절)};
\end{scope}

\end{tikzpicture}
\end{document}
